% chapter1-body.tex — pure content, no preamble.
% This file is \input{} by both main.tex (book) and chapter1.tex (article).




\section{Experiment}
I write a code in RUST to test how harmonic series
converge.

Here is what i found and I doubt.

In rust there are f32 \& f64 precison

With the simple code \cite[naive-method]{naive-method},
I can only finish the f32 with the following result.

To explain this I recall the definition of the f32
standars by IEEE, there are 23 digits of number and
8 number of exponent and 1 digits in sign.

We need find \(\frac{1}{n}\le \frac{1}{2}\Delta\)

Which the \(\Delta\) is the step size of the specific
range.

With this case is \(H_{2^{21}}\approx 15\),
and the \(\Delta\) in this range is \(2^{-23}*2^{3}=2^{-20}\)

According to IEEE standars when we hit \(2^{21}\)
we exactly have \(\frac{1}{2^{21}}=\frac{1}{2}\Delta\).

Then no more step are made, our program terminate.

But according to the same logic we can derive,
when the precision to f64, we have to stop at rough,
\(H_{2^{48}}\approx 34\),

With the step size \(\Delta=2^{-47}\),
\(\frac{1}{2^{48}}=\frac{1}{2}\Delta\).

But this would take about \(10^{6}s\) which is a really
long time.

With the same idea when \(\frac{1}{n}\le \Delta\) but
\(\frac{1}{n}\ge \frac{1}{2}\Delta\).
A lot term are bigger and simply equal \(\Delta\).
So with help with ChatGPT
I have \cite[accelerated]{accelerated}, to fast get the
result.

And verifies this idea within f32 \cite[verifyf32]{verifyf32}

Here is my Results.
\begin{lstlisting}
Numerically stopped.
n = 2097152
sum = 15.403683
1/n = 0.00000047683716
\end{lstlisting}

\begin{lstlisting}
accelerator starts:
  first n to process = 1675
  current sum        = 8.000487327575684
  bits               = 01000001000000000000000111111111

naive:
  n     = 2097152
  sum   = 15.403682708740234
  1/n   = 4.768371582031250e-7
  bits  = 01000001011101100111010101111100

accelerated:
  n     = 2097152
  sum   = 15.403682708740234
  1/n   = 4.768371582031250e-7
  bits  = 01000001011101100111010101111100
  blocks= 1252

PASS: naive and accelerated f32 end bit-for-bit identically.
\end{lstlisting}

\begin{lstlisting}
entered [16, 32):
last n = 4989191
sum    = 16.00000009545252055

Numerically stopped.
n   = 281474976710657
sum = 34.12203566804787158
1/n = 3.55271367880048831e-15
\end{lstlisting}

Professor I have analyzed your slides' numericals
it is , \(10^{15}\) with \(s_n=35.116\) which I guess your
precision system is some what a F64 systme.

and for the \(\frac{1}{k^2}\) problem is
\(n=94,906,266\) whic is highly close to \(2^{53/2}\)

I guess you also use a f64 system For the same logic as follow.

\(\zeta(2)\approx 1.64\) with this range the \(\Delta=2^{-52}\)

we need \(\frac{1}{k^2}\le \frac{1}{2}\Delta\)

which is \(k\approx 2^{\frac{53}{2}}\)

Finally, As I already guressed your numerical system digits,
I have a chanllenge question
to ask you, This two result the square is okay to
wait, but how you get the result for the harmonic series,
with my computer have \(\frac{10^9tick}{second}\) i need to run
\(10^{15-9}=10^{6}s\) which is a very long time(11.5 days)?

I do this by the code \cite[accelerated]{accelerated},
but how you had done to get the result?


% \subsection{本文问题定义}
% 对于$6\mid abc中(a,b,c)=(2,5,3m)$这一族，尝试证明CKW成立。
% 也就是
% $$
%     \boxed{
%         p_a,p_b,p_c\text{ 是正则序列}
%         \Leftarrow
%         (a,b,c)=(2,5,3m),\forall m\geq 1.
%     }
% $$


% \subsection{证明过程}
% Step1

% 我们要证明

% $$
%     p_2,p_3,p_5
% $$

% regular。

% 也就是证明

% $$
%     \begin{cases}
%         x^2+y^2+z^2=0, \\
%         x^3+y^3+z^3=0, \\
%         x^5+y^5+z^5=0
%     \end{cases}
% $$

% 只能推出

% $$
%     x=y=z=0.
% $$

% 全部换成基本对称多项式

% $$
%     e_1,e_2,e_3.
% $$

% Step2

% 由 Newton 恒等式\cite{NId}：

% $$
%     p_1=e_1.
% $$

% 以及

% $$
%     p_2=e_1p_1-2e_2.
% $$

% 因为 $p_1=e_1$，所以

% $$
%     \boxed{p_2=e_1^2-2e_2}.
% $$

% 现在假设

% $$
%     p_2=0.
% $$

% 那么

% $$
%     e_1^2-2e_2=0,
% $$

% 所以

% $$
%     \boxed{e_2=\frac{e_1^2}{2}}.
% $$


% 原来三个未知量

% $$
%     (e_1,e_2,e_3)
% $$

% 现在 $e_2$ 已经由 $e_1$ 决定了。

% 只剩下

% $$
%     e_1,e_3.
% $$

% Step3

% 加入 $p_3=0$

% Newton 恒等式给

% $$
%     p_3=e_1p_2-e_2p_1+3e_3.
% $$

% 因为

% $$
%     p_2=0,\qquad p_1=e_1,
% $$

% 所以

% $$
%     p_3=-e_1e_2+3e_3.
% $$

% 要求

% $$
%     p_3=0,
% $$

% 于是

% $$
%     3e_3=e_1e_2.
% $$

% 刚才已经知道

% $$
%     e_2=\frac{e_1^2}{2}.
% $$

% 代入：

% $$
%     3e_3
%     =
%     e_1\frac{e_1^2}{2}
%     =
%     \frac{e_1^3}{2}.
% $$

% 所以

% $$
%     \boxed{e_3=\frac{e_1^3}{6}}.
% $$


% $$
%     \boxed{
%         p_2=p_3=0
%     }
% $$

% 已经强迫

% $$
%     \boxed{
%         e_2=\frac12e_1^2,\qquad
%         e_3=\frac16e_1^3.
%     }
% $$

% 所以原来三维的 $(e_1,e_2,e_3)$ 空间，现在只剩下一个自由参数：

% $$
%     e_1.
% $$

% Step4

% 现在看 $p_5$

% 我们甚至不需要展开 $p_5$。

% Newton 递推是

% $$
%     p_k=e_1p_{k-1}-e_2p_{k-2}+e_3p_{k-3}.
% $$

% 先取 $k=4$：

% $$
%     p_4=e_1p_3-e_2p_2+e_3p_1.
% $$

% 而我们已经假定

% $$
%     p_2=p_3=0.
% $$

% 所以前两项消失：

% $$
%     p_4=e_3p_1.
% $$

% 又因为

% $$
%     p_1=e_1,
% $$

% 所以

% $$
%     \boxed{p_4=e_1e_3}.
% $$

% 然后取 $k=5$：

% $$
%     p_5=e_1p_4-e_2p_3+e_3p_2.
% $$

% 后两项又因为

% $$
%     p_2=p_3=0
% $$

% 而消失。

% 所以

% $$
%     p_5=e_1p_4.
% $$

% 而 $p_4=e_1e_3$，因此

% $$
%     \boxed{p_5=e_1^2e_3}.
% $$

% 再代入

% $$
%     e_3=\frac{e_1^3}{6},
% $$

% 得到

% $$
%     \boxed{p_5=\frac{e_1^5}{6}}.
% $$

% 现在第三个方程要求

% $$
%     p_5=0.
% $$

% 因此

% $$
%     \frac{e_1^5}{6}=0,
% $$

% 只能有

% $$
%     \boxed{e_1=0}.
% $$

% 于是前面的公式立刻给：

% $$
%     e_2=\frac{e_1^2}{2}=0,
% $$

% $$
%     e_3=\frac{e_1^3}{6}=0.
% $$

% 所以

% $$
%     \boxed{e_1=e_2=e_3=0}.
% $$


% LastStep

% 使用维达

% $x,y,z$ 是

% $$
%     T^3-e_1T^2+e_2T-e_3
% $$

% 的三个根。

% 现在

% $$
%     e_1=e_2=e_3=0,
% $$

% 所以这个多项式就是

% $$
%     T^3.
% $$

% 它唯一的根是

% $$
%     0.
% $$

% 因此

% $$
%     \boxed{x=y=z=0}.
% $$

% 是一个三重根，是唯一的

% 所以我们证明了

% $$
%     \boxed{p_2,p_3,p_5\text{ regular}.}
% $$




% \section{英文证明}


% % \subsection{MVHR}
% % 本问题理论部分可以参考\cite[MVHR]{MVHR}模型\\
% % 先把我们的问题严格写出来
% % 假设我们现在有：
% % $$
% %     Q_S=100\text{ 吨豆粕}
% % $$
% % 现货价格变化：
% % $$
% %     \Delta S
% % $$

% % 期货价格变化：

% % $$
% %     \Delta F
% % $$

% % 我们决定 hedge 比例：

% % $$
% %     h
% % $$

% % 那么每吨对应的套保 PnL 变化可以写成：

% % $$
% %     \Delta P=\Delta S-h\Delta F
% % $$

% % 我们的目标非常朴素：

% % 找一个 $h$，让这个组合尽可能不波动。

% % 也就是：

% % $$
% %     \min_h Var(\Delta S-h\Delta F)
% % $$

% % 把它展开：

% % $$
% %     Var(\Delta S-h\Delta F)
% %     =
% %     Var(\Delta S)
% %     +h^2Var(\Delta F)
% %     -2hCov(\Delta S,\Delta F)
% % $$

% % 对 $h$ 求导：

% % $$
% %     2hVar(\Delta F)
% %     -
% %     2Cov(\Delta S,\Delta F)
% %     =0
% % $$

% % 于是：

% % $$
% %     \boxed{
% %         h^*
% %         =
% %         \frac{Cov(\Delta S,\Delta F)}
% %         {Var(\Delta F)}
% %     }
% % $$

% % 这就是我们现在要用的核心公式。


% % \subsection{数据获取}
% % 数据从网站项目获取\cite[akshare]{akshare}，此软件的文档
% % 在\cite[doc]{akshare-doc}。
% % 数据获取实验在\textbf{MVHR-akshare}分支。




\section{Reference}
\bibliographystyle{plain}  % or another style like unsrt, IEEEtran, etc.
\bibliography{chapter1-body-references}  % references.bib is the file name


% \subsection{Definition of EMH}
% \begin{definition}[EMH means Symmetries]
%     \label{def:mydefinition}
%     We say a market is completely EMH if and
%     only if for every condition $A$, the random
%     variable of the logarithmic growth rate $W$
%     agrees that $\mathbb{E}_{W\mid A}(W)\doteq 0$.
% \end{definition}


% \subsection{Prior Work}
% Discussion of what others have done.




% \section{Methods}
% How you approached the problem.




% \section{Results}
% What you found.




% \section{Conclusion}
% Your closing thoughts.




